% =====================================================================
%  GMN2 附录 C  WKB 误差分析（WKB error analysis）
%  原文：wkb-paper.tex 9673–9757 行；BRIEF 印刷页 169–171（ch19 块，整块 28 式）；编号公式 (C.1)–(C.11) 共 11 条，另有脚注 1 条
%  本块引用全在本文件内，\eqref 可直接用；\Pexp→\operatorname{Pexp}、\id→\mathrm{id}（preamble 未定义）
% =====================================================================

\section{WKB 误差分析}
\label{app:wkb-error}

本附录考虑指数式增长的平坦截面沿 WKB 曲线的传播。我们想看清的是：这种传播在 $\zeta\to0$ 时的渐近行为，恰好就是把 Higgs 场 $\varphi$ 的本征值 $\lambda$ 积分出来所得的结果。

更精确地说：取一个使 $\varphi$ 对角化的规范，
\begin{equation}
\varphi = \begin{pmatrix} \lambda & 0 \\ 0 & -\lambda \end{pmatrix}.
\end{equation}
然后令 $z(t)$ 是一条相位为 $\vartheta$ 的 WKB 曲线，并取 $\zeta\in\IH_\vartheta$，使得 $\re\lambda_z z'(t)/\zeta<0$。设 $s$ 是一个平坦截面，满足
\begin{equation}
s(z(0)) = \begin{pmatrix} 1 \\ 0 \end{pmatrix}.
\end{equation}
WKB 近似的命题是：当 $\zeta\to0$ 时
\begin{equation} \label{eq:wkb}
s(z(t)) \sim c(t) \begin{pmatrix} e^{-\frac{R}{\zeta} \int_{z(0)}^{z(t)} \lambda} \\ 0 \end{pmatrix},
\end{equation}
其中 $c(t)$ 是某个与 $\zeta$ 无关的函数。

为证明 \eqref{eq:wkb}，先如下定义 WKB 余项 $\psi$：
\begin{equation} \label{eq:wkb-remainder}
\psi(z) = s(z) \exp\left( \frac{R}{\zeta} \int_{z(0)}^z \lambda \right).
\end{equation}
平坦性方程 $(d+\CA)s=0$ 于是变成
\begin{align}
(\partial_{\bar z} + A_{\bar z} + R\zeta \bar\varphi_{\bar z}) \psi &= 0, \\
(\partial_z + A_z + \frac{R}{\zeta} (\varphi_z - \lambda_z\, \mathrm{id})) \psi &= 0.
\end{align}
所以沿曲线 $z(t)$，$\psi$ 的演化是
\begin{equation} \label{eq:reduced-ode}
\left( \frac{d}{dt} + B(t) \right) \psi(z(t)) = 0
\end{equation}
其中
\begin{equation}
B(t) = \frac{R}{\zeta} z' \begin{pmatrix} 0 & 0 \\ 0 & -2\lambda_z \end{pmatrix} + z' A_z + \bar z' (A_{\bar z} + R\zeta \bar\varphi_{\bar z}).
\end{equation}

所要的 \eqref{eq:wkb} 等价于说：$\lim_{\zeta\to0}\psi$ 存在，而且其第二个分量消失。为什么应当如此？直觉是：$B(t)$ 中当 $\zeta\to0$ 时不为有限的项，只有右下角那一项 $-2Rz'\lambda_z/\zeta$；而这一项的唯一效果，是在 $\psi$ 的第二个分量中引入一个形如 $e^{-\#/\zeta}$ 的因子（其中 $\re\# > 0$），于是该分量在 $\zeta\to0$ 时消失。

为论证这一直觉，令 $B_0(t)=g(t)B(t)g^{-1}(t)$ 是 $B(t)$ 的对角化。我们希望证明
\begin{equation} \label{eq:wkb-interm}
\lim_{\zeta\to0} \bigg\lVert \operatorname{Pexp}\int_0^t B(t) - \operatorname{Pexp}\int_0^t B_0(t) \bigg\rVert = 0.
\end{equation}
一旦确立了 \eqref{eq:wkb-interm}，就可以在 $\zeta\to0$ 极限下用 $B_0$ 代替 $B$ 来计算 $\psi$ 的平行移动；那便证明了所要的 \eqref{eq:wkb}。

一个直接的计算——只用到一个事实：$B(t)$ 是 $2\times2$ 矩阵，其右下角元素的实部趋于 $-\infty$，而其余元素都保持有限——表明 $\exp B_0(t)$ 在 $\zeta\to0$ 时有界，并且可以选取 $g(t)$ 使 $g(t)\to\mathrm{id}$（当 $\zeta\to0$ 时）。由这两件事即得 \eqref{eq:wkb-interm}。为看清这一点，我们先把区间切成小段，使 $B$ 在每段上缓变，从而归结为 $B$ 与 $t$ 无关的情形。此时 \eqref{eq:wkb-interm} 化为
\begin{equation}
\lim_{\zeta\to0} \lVert e^{tB} - e^{tB_0} \rVert = 0.
\end{equation}
而这又可由 $B=gB_0g^{-1}$、$\lim_{\zeta\to0}g=\mathrm{id}$ 以及 $\lim_{\zeta\to0}e^{tB_0}$ 的存在性推出。\footnote{正是在这里，我们用到了如下事实：当 $\zeta\to0$ 时，$B(t)$ 右下角元素的实部趋于 $-\infty$，而非 $+\infty$。} 这就完成了 \eqref{eq:wkb-interm} 从而 \eqref{eq:wkb} 的证明。

更精确地说，我们是在 WKB 曲线始于正则点的情形下证明 \eqref{eq:wkb} 的。而我们的应用需要考虑 WKB 曲线始于奇点的情形，不妨设该奇点位于 $z(0)$。这时显然不能指望 \eqref{eq:wkb} 原封不动地成立，因为 $\int_{z(0)}^{z(t)}\lambda$ 发散。替代的办法是另取一个函数 $I(z)$，它定义在 WKB 曲线的一个邻域上（奇点除外），满足 $dI=\lambda$。于是与上面类似的论证表明，存在一个平坦截面 $s$，它在 $\zeta\to0$ 时的行为是
\begin{equation} \label{eq:wkb-sing}
s(z(t)) \sim \begin{pmatrix} c(t) e^{-\frac{R}{\zeta} I(z(t))} \\ 0 \end{pmatrix}.
\end{equation}
而且可以取 $s$ 为与 WKB 曲线 $z(t)$ 及奇点 $z(0)$ 相伴的那个小平坦截面。
